%% 第二章 相关理论与关键技术
\chapter{相关理论与关键技术}

本章介绍本文系统所依赖的理论与技术基础，包括目标检测与卷积神经网络、YOLOv8 的网络结构、注意力机制、卡尔曼滤波、针孔相机模型、PX4 与 jMAVSim 仿真框架，以及 PySide6 与网络通信等内容。理解这些基础有助于把握后续章节中各项设计选择的依据。

\section{目标检测基础}

\subsection{卷积神经网络}

卷积神经网络（Convolutional Neural Network，CNN）通过局部连接、权值共享与空间池化三种机制，从图像中逐层提取由低级到高级的特征。设第 $l$ 层的输入特征图为 $\mathbf{X}^{l-1}\in\mathbb{R}^{H_{l-1}\times W_{l-1}\times C_{l-1}}$，卷积核为 $\mathbf{W}^{l}\in\mathbb{R}^{k\times k\times C_{l-1}\times C_l}$，则第 $l$ 层的输出为

\begin{equation}
\mathbf{X}^{l}_{i,j,c} = \sigma\left(\sum_{u=0}^{k-1}\sum_{v=0}^{k-1}\sum_{c'=0}^{C_{l-1}-1}
\mathbf{W}^{l}_{u,v,c',c}\,\mathbf{X}^{l-1}_{i+u,\,j+v,\,c'} + b^{l}_{c}\right),
\label{eq:conv}
\end{equation}

其中 $k$ 为卷积核尺寸，$b^l_c$ 为偏置，$\sigma(\cdot)$ 为非线性激活函数。现代检测网络普遍采用 SiLU（Swish）激活：

\begin{equation}
\mathrm{SiLU}(x) = x\cdot \mathrm{sigmoid}(x) = \frac{x}{1+e^{-x}}.
\label{eq:silu}
\end{equation}

与 ReLU 相比，SiLU 在负半轴保留小幅非零梯度，有助于缓解小目标训练中的梯度消失。

批归一化（Batch Normalization，BN）在每个小批量内对通道维度做标准化，使各层输入分布保持稳定，从而允许更大的学习率与更快的收敛：

\begin{equation}
\hat{x} = \frac{x-\mu_{\mathcal{B}}}{\sqrt{\sigma_{\mathcal{B}}^2+\epsilon}},\qquad
y = \gamma\hat{x}+\beta.
\label{eq:bn}
\end{equation}

\subsection{特征金字塔与多尺度检测}

目标在图像中的尺度差异极大，单一尺度的特征图难以同时兼顾大小目标。特征金字塔网络（Feature Pyramid Network，FPN）\cite{fpn} 通过自顶向下的路径将深层语义特征与浅层高分辨率特征融合，使每一层都同时具备较强的语义与空间信息。PANet\cite{panet} 在 FPN 基础上增加自底向上的路径，进一步缩短浅层特征到检测头的传播距离。YOLOv8 的颈部即采用 PAN-FPN 结构。

设骨干网络输出的特征图为 $C_3,C_4,C_5$，分别对应 stride 8、16、32；FPN 的融合过程可写为

\begin{equation}
P_l = \mathrm{Conv}\left(\mathrm{Upsample}(P_{l+1}) \oplus \mathrm{Conv}(C_l)\right),
\label{eq:fpn}
\end{equation}

其中 $\oplus$ 表示逐元素相加，$\mathrm{Conv}$ 通常为 $1\times1$ 卷积（用于通道对齐）加 $3\times3$ 卷积（用于消除上采样混叠）。PAN 路径则从 $P_3$ 向 $P_5$ 做下采样融合。最终检测头分别接在 $P_3,P_4,P_5$ 上，对应的下采样倍数越大，感受野越大，适合检测大目标；倍数越小，空间分辨率越高，适合检测小目标。

\subsection{检测头的预测形式与损失函数}

YOLOv8 采用无锚框（anchor-free）的检测头，每个网格位置直接预测目标中心相对网格的偏移量、宽高以及类别概率。对于分类分支，使用二元交叉熵（BCE）损失：

\begin{equation}
L_{cls} = -\frac{1}{N_{pos}}\sum_{i}\left[y_i\log p_i + (1-y_i)\log(1-p_i)\right].
\label{eq:bce}
\end{equation}

对于边框回归，YOLOv8 使用分布焦点损失（Distribution Focal Loss，DFL）与 CIoU 损失的组合。DFL 将边框回归建模为离散概率分布，使网络能够学习边界的不确定性：

\begin{equation}
L_{DFL} = -\sum_{i}\left((y_{i+1}-y)\log p_i + (y-y_i)\log p_{i+1}\right).
\label{eq:dfl}
\end{equation}

CIoU 同时考虑重叠面积、中心点距离与宽高比一致性：

\begin{equation}
L_{CIoU} = 1 - \mathrm{IoU} + \frac{\rho^2(\mathbf{b},\mathbf{b}^{gt})}{c^2} + \alpha v,
\label{eq:ciou}
\end{equation}

其中 $\rho(\cdot)$ 为欧氏距离，$c$ 为最小外接矩形对角线长度，$v$ 为宽高比一致性项，$\alpha$ 为权重系数。

\subsection{评测指标}

目标检测的常用指标包括 IoU、Precision、Recall、AP 与 mAP。交并比（IoU）定义为预测框与真值框交集面积与并集面积之比：

\begin{equation}
\mathrm{IoU} = \frac{|\mathcal{B}_{pred}\cap\mathcal{B}_{gt}|}{|\mathcal{B}_{pred}\cup\mathcal{B}_{gt}|}.
\label{eq:iou}
\end{equation}

在某一 IoU 阈值下，将预测按置信度降序排列，逐点计算 Precision 与 Recall，可得到 P-R 曲线，其下面积即为该类别在该阈值下的平均精度 AP。对所有类别取平均得到 mAP。COCO 评测协议以 IoU 从 0.50 到 0.95、步长 0.05 的十个阈值取平均，记为 mAP50-95；mAP50 则只取 IoU=0.50。

COCO 还按目标框面积将目标划分为小、中、大三类：面积小于 $32^2$ 像素为 small，介于 $32^2$ 与 $96^2$ 之间为 medium，大于 $96^2$ 为 large。对于反无人机任务，small 子集上的 AP（本文记为 AP\_small）是最具参考价值的指标，因为远距离目标大多落在该区间。需要强调的是，检测框面积与目标实际物理尺寸无关，同一架无人机在不同距离下会落入不同的分桶，因此 AP\_small 反映的是"模型在多远距离上仍能检出"的综合能力。

速度指标通常以 FPS（Frames Per Second）衡量。由于不同实现的预处理、后处理与显存拷贝开销差异很大，FPS 必须注明测量条件。本文统一采用 batch=1、帧预载内存、包含预处理与后处理的口径，以保证不同模型之间的可比性。

\section{YOLOv8 网络结构}

\subsection{整体结构}

YOLOv8 由骨干（Backbone）、颈部（Neck）与头部（Head）三部分组成。骨干采用改进的 CSPDarknet，通过 C2f 模块在保持梯度流动的同时减少计算量；颈部采用 PAN-FPN，实现多尺度特征的双向融合；头部采用解耦结构，将分类与回归分支分离，避免两个任务的特征相互干扰。YOLOv8n 是五种规模中最小的一种，参数量约 3.0\,M，在 640$\times$640 输入下的计算量约 8.7\,GFLOPs。

\subsection{C2f 模块}

C2f（CSP Bottleneck with 2 convolutions and f fused outputs）模块是 YOLOv8 的核心组件。其结构可以概括为：输入先经过 $1\times1$ 卷积降维，然后分为两支，一支直接短路，另一支经过若干个 Bottleneck 后逐级拼接，最后再经过 $1\times1$ 卷积融合。与 YOLOv5 的 C3 模块相比，C2f 保留了更多的梯度流分支，使浅层特征能够更直接地传递到检测头，这对小目标检测尤为重要。

设输入为 $\mathbf{X}$，C2f 的输出可写为

\begin{equation}
\mathbf{Y} = \mathrm{Conv}_{1\times1}\left(\mathrm{Concat}\left[\mathbf{X}_0,
\mathcal{B}_1(\mathbf{X}_1),\mathcal{B}_2(\mathbf{X}_2),\ldots,\mathcal{B}_n(\mathbf{X}_n)\right]\right),
\label{eq:c2f}
\end{equation}

其中 $\mathcal{B}_i$ 表示第 $i$ 个 Bottleneck，$\mathbf{X}_i$ 为前一层输出。由于每个 Bottleneck 的输出都参与最终拼接，网络实际上获得了多感受野的特征组合。

\subsection{SPPF 与解耦头}

SPPF（Spatial Pyramid Pooling - Fast）通过串联多个不同尺寸的最大池化层，在不显著增加计算量的前提下扩大感受野，增强模型对目标整体结构的感知。解耦头则将分类与回归分别送入两条独立的卷积分支，分类分支输出 $C$ 个类别的概率，回归分支输出 4 个边框参数与 16 个分布参数。解耦设计减少了任务间的负迁移，是 YOLOv8 相比前代的重要改进之一。

\section{注意力机制}

\subsection{从 SE 到 CBAM}

Squeeze-and-Excitation（SE）模块通过全局平均池化获得通道描述子，再经两层全连接网络生成通道权重：

\begin{equation}
\mathbf{s} = \sigma\left(\mathbf{W}_2\,\delta\left(\mathbf{W}_1\,\mathrm{GAP}(\mathbf{X})\right)\right),\qquad
\mathbf{X}' = \mathbf{s}\odot\mathbf{X}.
\label{eq:se}
\end{equation}

CBAM\cite{cbam} 在通道注意力之后串联空间注意力，空间分支对通道维度做最大池化与平均池化，再经 $7\times7$ 卷积生成空间权重图。二者的共同点是"先压缩、再激励"，但都只在一个维度上建模，且空间分支的感受野有限。

\subsection{EMA 高效多尺度注意力}

EMA（Efficient Multi-scale Attention）\cite{ema} 的核心思想是\textbf{分组并行 + 跨空间学习}。设输入特征图为 $\mathbf{X}\in\mathbb{R}^{C\times H\times W}$，EMA 首先沿通道维度将其分为 $G$ 组，每组独立进行如下处理：

（1）对组内特征分别沿水平与垂直方向做一维池化，得到两组方向感知的描述子；

（2）对描述子做一维卷积，建模局部跨通道交互；

（3）将两组描述子通过 sigmoid 归一化后与原始特征逐元素相乘，得到方向注意力增强的特征；

（4）对增强后的特征做全局平均池化，并通过 softmax 生成跨空间的注意力权重；

（5）将两组权重与原始输入相乘后相加，输出最终特征。

与 CBAM 相比，EMA 的优势在于：分组并行降低了大通道数下的计算量；一维方向卷积以极小的参数量建模长程依赖；跨空间学习使不同组之间能够交换信息，避免分组造成的特征割裂。对于小目标而言，这种"方向感知 + 跨组交互"的机制能够在复杂背景中更准确地定位目标所在的条带区域。

\section{卡尔曼滤波}

\subsection{线性高斯状态空间模型}

卡尔曼滤波\cite{kalman} 解决的是线性高斯系统的最优状态估计问题。设离散时间系统的状态方程与量测方程为

\begin{equation}
\mathbf{x}_k = \mathbf{F}_k\mathbf{x}_{k-1} + \mathbf{B}_k\mathbf{u}_k + \mathbf{w}_k,\qquad
\mathbf{z}_k = \mathbf{H}_k\mathbf{x}_k + \mathbf{v}_k,
\label{eq:state_space}
\end{equation}

其中 $\mathbf{w}_k\sim\mathcal{N}(\mathbf{0},\mathbf{Q}_k)$ 为过程噪声，$\mathbf{v}_k\sim\mathcal{N}(\mathbf{0},\mathbf{R}_k)$ 为量测噪声，二者相互独立。卡尔曼滤波分为预测与更新两步。

预测步：

\begin{equation}
\hat{\mathbf{x}}_{k|k-1} = \mathbf{F}_k\hat{\mathbf{x}}_{k-1|k-1},
\label{eq:kf_predict_x}
\end{equation}

\begin{equation}
\mathbf{P}_{k|k-1} = \mathbf{F}_k\mathbf{P}_{k-1|k-1}\mathbf{F}_k^{T} + \mathbf{Q}_k.
\label{eq:kf_predict_P}
\end{equation}

更新步：

\begin{equation}
\mathbf{K}_k = \mathbf{P}_{k|k-1}\mathbf{H}_k^{T}\left(\mathbf{H}_k\mathbf{P}_{k|k-1}\mathbf{H}_k^{T}+\mathbf{R}_k\right)^{-1},
\label{eq:kf_gain}
\end{equation}

\begin{equation}
\hat{\mathbf{x}}_{k|k} = \hat{\mathbf{x}}_{k|k-1} + \mathbf{K}_k\left(\mathbf{z}_k-\mathbf{H}_k\hat{\mathbf{x}}_{k|k-1}\right),
\label{eq:kf_update}
\end{equation}

\begin{equation}
\mathbf{P}_{k|k} = \left(\mathbf{I}-\mathbf{K}_k\mathbf{H}_k\right)\mathbf{P}_{k|k-1}.
\label{eq:kf_cov}
\end{equation}

增益矩阵 $\mathbf{K}_k$ 的物理含义是"在预测与量测之间分配信任"：当量测噪声 $\mathbf{R}$ 较大时，增益变小，滤波器更信任预测；当预测不确定度 $\mathbf{P}$ 较大时，增益变大，滤波器更信任量测。

\subsection{自适应噪声估计}

标准卡尔曼滤波要求 $\mathbf{Q}$ 与 $\mathbf{R}$ 已知且恒定。在视觉跟踪中，检测框的抖动程度随目标距离、光照与背景而变化，固定 $\mathbf{R}$ 难以适应。Sage-Husa 自适应滤波\cite{sagehusa} 利用新息（innovation）序列在线估计噪声统计特性。新息定义为

\begin{equation}
\boldsymbol{\nu}_k = \mathbf{z}_k - \mathbf{H}_k\hat{\mathbf{x}}_{k|k-1}.
\label{eq:innovation}
\end{equation}

在理论噪声统计下，新息应为零均值白噪声，其协方差为

\begin{equation}
\mathbf{S}_k = \mathbf{H}_k\mathbf{P}_{k|k-1}\mathbf{H}_k^{T}+\mathbf{R}_k.
\label{eq:innovation_cov}
\end{equation}

通过对新息协方差做滑动平均，可以反解出 $\mathbf{R}_k$ 的估计：

\begin{equation}
\hat{\mathbf{R}}_k = \frac{1}{M}\sum_{i=k-M+1}^{k}\left(\boldsymbol{\nu}_i\boldsymbol{\nu}_i^{T}
-\mathbf{H}_i\mathbf{P}_{i|i-1}\mathbf{H}_i^{T}\right),
\label{eq:sage_husa}
\end{equation}

其中 $M$ 为滑动窗口长度。实际实现中通常对 $\hat{\mathbf{R}}_k$ 做下界约束与对角化处理，以避免数值不稳定。本文在第六章给出具体的实现细节。

\subsection{滤波的收敛性与可观测性}

卡尔曼滤波的估计误差协方差在系统一致可观测的条件下收敛。对于本文的六维状态（位置与速度），量测只包含位置，速度通过位置的时间差分间接观测。若目标长时间匀速直线运动，速度的可观测性较弱，滤波收敛较慢；若目标做机动，位置的变化会提供额外的速度信息。工程上通常通过增大过程噪声 $\mathbf{Q}$ 中的速度项来加快收敛，但代价是稳态抖动增大。本文在参数整定中对二者做了折中。

\section{针孔相机模型与标定}

\subsection{坐标系定义}

本文涉及四个坐标系：世界坐标系、相机坐标系、图像坐标系与像素坐标系。世界坐标系采用北—东—地（North-East-Down，NED）约定，$x$ 轴指向北，$y$ 轴指向东，$z$ 轴指向地心，地面上 $z=0$，高度越大 $z$ 越小。相机坐标系以光心为原点，$z_c$ 轴沿光轴指向景物，$x_c$ 轴向右，$y_c$ 轴向下。图像坐标系以主点为原点、以物理尺寸为单位；像素坐标系以图像左上角为原点、以像素为单位。

设世界坐标系到相机坐标系的旋转矩阵为 $\mathbf{R}$、平移向量为 $\mathbf{t}$，则点 $P_w$ 在相机坐标系下的坐标为

\begin{equation}
\mathbf{P}_c = \mathbf{R}\,\mathbf{P}_w + \mathbf{t}.
\label{eq:world2cam}
\end{equation}

\subsection{投影模型}

针孔模型将三维点投影到归一化像平面，再经内参矩阵映射到像素坐标：

\begin{equation}
s\begin{bmatrix}u\\v\\1\end{bmatrix}
= \mathbf{K}\begin{bmatrix}\mathbf{R}&\mathbf{t}\end{bmatrix}
\begin{bmatrix}X_w\\Y_w\\Z_w\\1\end{bmatrix},\qquad
\mathbf{K}=\begin{bmatrix}f_x&0&c_x\\0&f_y&c_y\\0&0&1\end{bmatrix},
\label{eq:projection}
\end{equation}

其中 $s=Z_c$ 为深度，$f_x,f_y$ 为以像素为单位的焦距，$(c_x,c_y)$ 为主点。若只考虑相机坐标系下的点，则

\begin{equation}
u = f_x\frac{X_c}{Z_c}+c_x,\qquad v = f_y\frac{Y_c}{Z_c}+c_y.
\label{eq:pixel}
\end{equation}

实际镜头存在径向畸变与切向畸变，常用的 Brown-Conrady 模型为

\begin{equation}
\begin{aligned}
x_d &= x(1+k_1r^2+k_2r^4+k_3r^6)+2p_1xy+p_2(r^2+2x^2),\\
y_d &= y(1+k_1r^2+k_2r^4+k_3r^6)+p_1(r^2+2y^2)+2p_2xy,
\end{aligned}
\label{eq:distortion}
\end{equation}

其中 $(x,y)$ 为归一化无畸变坐标，$r^2=x^2+y^2$，$k_1,k_2,k_3$ 为径向畸变系数，$p_1,p_2$ 为切向畸变系数。本文仿真相机无畸变，实拍相机在标定后按上述模型去畸变。

\subsection{张正友标定法}

张正友标定法\cite{zhang_calib} 利用平面棋盘格在不同位姿下的多幅图像，通过单应性矩阵约束求解内参。设棋盘格平面为 $Z_w=0$，则式~\eqref{eq:projection} 退化为单应变换：

\begin{equation}
s\begin{bmatrix}u\\v\\1\end{bmatrix}
= \mathbf{K}\begin{bmatrix}\mathbf{r}_1&\mathbf{r}_2&\mathbf{t}\end{bmatrix}
\begin{bmatrix}X_w\\Y_w\\1\end{bmatrix}
= \mathbf{H}\begin{bmatrix}X_w\\Y_w\\1\end{bmatrix}.
\label{eq:homography}
\end{equation}

利用旋转向量正交约束 $\mathbf{r}_1^{T}\mathbf{r}_2=0$、$\|\mathbf{r}_1\|=\|\mathbf{r}_2\|=1$，可以构造关于 $\mathbf{K}$ 的线性方程组，至少三幅不同位姿的图像即可求解。随后用最大似然估计优化重投影误差：

\begin{equation}
\min_{\mathbf{K},\mathbf{k},\{\mathbf{R}_i,\mathbf{t}_i\}}\sum_{i}\sum_{j}
\left\|\mathbf{p}_{ij}-\hat{\mathbf{p}}(\mathbf{K},\mathbf{k},\mathbf{R}_i,\mathbf{t}_i,\mathbf{P}_j)\right\|^2.
\label{eq:reproj}
\end{equation}

本文在自建数据采集环节使用 9$\times$6、方格边长 25\,mm 的棋盘格完成内参标定，标定结果用于评估单目定位的绝对精度。

\section{PX4 与 jMAVSim 仿真框架}

\subsection{PX4 Autopilot 与 SITL}

PX4\cite{px4} 是一款开源飞控固件，支持多旋翼、固定翼、垂直起降等多种机型。其软件在环（SITL）模式将飞控固件编译为可在普通 PC 上运行的进程，通过 MAVLink\cite{mavlink} 协议与仿真器交换传感器数据与执行器指令。SITL 的优势在于：无需真实飞控硬件即可验证控制逻辑、航点任务与失效保护；可以与多种仿真器（jMAVSim、Gazebo、AirSim）组合；支持 lockstep 模式，使仿真时间与飞控时间严格同步，保证实验结果可复现。

在 lockstep 模式下，仿真器每推进一步都等待飞控完成相应的控制计算，二者以相同的仿真时间前进。这使得"仿真 60 秒"与"飞控内部 60 秒"完全一致，避免了因计算速度差异导致的时间漂移。本文的所有闭环实验均启用 lockstep。

\subsection{jMAVSim 架构}

jMAVSim\cite{jmavsim} 是 PX4 自带的轻量级仿真器，使用 Java 与 Java3D 构建。其核心类包括：

（1）\texttt{Simulator}：仿真主循环，负责解析命令行参数、创建世界、启动物理步进与网络通信；

（2）\texttt{World}：物理世界容器，管理所有 \texttt{KinematicObject}，每个仿真步调用各对象的 \texttt{update} 方法；

（3）\texttt{KinematicObject}：所有可运动对象的基类，提供位置、速度、姿态接口；

（4）\texttt{Vehicle}：多旋翼动力学模型，接收 MAVLink 指令并积分运动方程；

（5）\texttt{Visualizer3D}：基于 Java3D 的三维渲染器，负责场景构建、相机视角与画面输出；

（6）\texttt{MAVLinkInterface}：与 PX4 的 MAVLink 通信接口。

jMAVSim 的坐标约定与 PX4 一致，采用 NED：$x$ 向北、$y$ 向东、$z$ 向下。场景中的地面位于 $z=0$，飞行器高度为负值。本文在改造 jMAVSim 时严格遵循这一约定，并在真值广播中加入坐标转换说明，避免与算法端的相机坐标系混淆。

\subsection{离屏渲染与帧流}

jMAVSim 默认依赖窗口系统显示三维画面。在无头服务器或 WSL 环境中，窗口系统可能不可用或不稳定，本文因此实现了纯离屏渲染模式：关闭 Java3D 的窗口画布，改用离屏 \texttt{Canvas3D} 渲染到缓冲区，再将帧编码为 JPEG 通过 TCP 推送。该模式避免了图形界面依赖，使仿真可以在后台长时间稳定运行。具体实现将在第七章详细介绍。

\section{系统开发技术}

\subsection{PySide6 与 QtWebEngine}

PySide6 是 Qt 6 的官方 Python 绑定，提供窗口、控件、绘图、定时器、网络与多媒体等完整能力。本文使用其 \texttt{QWidget} 体系构建主界面，使用 \texttt{QPainter} 自绘仪表与罗盘，使用 \texttt{QPixmap} 显示视频帧，使用 \texttt{QSettings} 的替代方案（JSON 文件）持久化配置。

地图部分使用 QtWebEngine 嵌入 Leaflet\cite{leaflet} 页面。Leaflet 是一个轻量级开源地图库，支持瓦片图层、标记、折线、多边形与自定义覆盖物，能够在本地 HTML 中加载 Esri 卫星影像。Python 端通过 \texttt{QWebChannel} 或 JavaScript 注入向页面推送目标位置，实现地图标记的实时更新。

\subsection{OpenCV 与 PyTorch}

OpenCV 负责图像解码、颜色空间转换、几何绘制与视频读写。本文使用 \texttt{cv2.VideoCapture} 读取离线视频与摄像头，使用 \texttt{cv2.imencode} 将帧编码为 JPEG，使用 \texttt{cv2.line}、\texttt{cv2.rectangle}、\texttt{cv2.putText} 绘制检测叠加层。

PyTorch 与 Ultralytics 负责模型加载与推理。YOLOv8 模型可直接通过 \texttt{YOLO("best.pt")} 加载，调用 \texttt{model.predict} 获得检测框、置信度与类别。为降低端到端延迟，本文在推理时关闭冗余的日志与可视化，并将输入尺寸固定为 640$\times$640。

\subsection{TCP 与 UDP 通信}

系统涉及两类网络通信。视频帧流使用 TCP：TCP 提供可靠、有序的字节流，适合体积较大的 JPEG 帧；协议为"一行文本头 + 二进制数据"，头部包含帧序号、时间戳、宽度、高度与字节数，接收端按字节数精确读取，避免粘包。真值广播与结果输出使用 UDP：UDP 无连接、开销小，适合高频小报文；本文的真值广播频率为 50\,Hz，结果输出频率与处理帧率一致。

需要说明的是，UDP 不保证送达，因此在结果输出中本文采用"最新值覆盖"策略：显示端只关心当前状态，丢失个别报文不会影响整体显示；同时通过序号字段检测丢包率，用于评估链路质量。

\section{本章小结}

本章介绍了本文所依赖的理论与技术基础。在检测方面，说明了卷积网络、特征金字塔、检测头预测形式与评测指标；在模型方面，剖析了 YOLOv8 的 C2f、SPPF 与解耦头结构；在注意力方面，对比了 SE、CBAM 与 EMA 的建模方式；在滤波方面，给出了卡尔曼滤波的预测—更新方程与 Sage-Husa 自适应噪声估计；在定位方面，建立了针孔相机模型与张正友标定法；在仿真方面，介绍了 PX4 SITL、jMAVSim 的类结构与 NED 坐标约定；在系统方面，概述了 PySide6、Leaflet、OpenCV、PyTorch 与 TCP/UDP 的使用方式。这些内容为后续章节的设计与实验提供了共同的语言基础。
